\section{Prescribing the constant vector by a Gram matrix}\label{sec:gram}

\subsection{The seed constants as residues}

At $n=1$ the only poles are $0$ and $-4$. This makes it possible
to express the rational constants by integrals along the vertical
line halfway between the poles. We first identify the rational
kernels that occur in these integrals.

To write alternating sums over odd integers compactly, let $\ch(k)$
be $0$ for even $k$, $1$ for $k\equiv1\pmod4$, and $-1$ for
$k\equiv3\pmod4$. Thus $\ch$ records the signs in the beta series.
It is periodic with period $4$ and multiplicative.

Let $T$ be a proper rational function with poles at $0,-4$, invariant
under $t\mapsto-4-t$, with partial fractions
\[
 T(t)=\sum_{i=1}^s\frac{a_i}{t^i}
         +\sum_{i=1}^s\frac{(-1)^ia_i}{(t+4)^i}.
\]
In $\sum_{k\ge1}\ch(k)T(k/b)$, shifting the second pole's sum
by $4b$ leaves the signs unchanged and removes its first $4b$
terms. The rational constant is therefore
\[
 q_b=-\sum_{i=1}^s(-1)^ia_ib^i
                           \sum_{k=1}^{4b}\frac{\ch(k)}{k^i}.
\]
Introduce the kernel
\[
 H_b(t)=\sum_{\substack{1\le k<4b\\k\text{ odd}}}
                         \frac{\ch(k)}{t+k/b}.
\]
The coefficient of $t^{i-1}$ in $(t+k/b)^{-1}$ is
$(-1)^{i-1}(b/k)^i$, so
\begin{equation}\label{eq:residue-constant}
 q_b=\Res_{t=0}T(t)H_b(t).
\end{equation}

On the midpoint line $t=-2+iy$, pair the terms indexed by $k$
and $4b-k$. Their distances from the midpoint have opposite signs,
as do their character weights. The result is a real rational
function of $y^2$:
\begin{equation}\label{eq:gb}
 H_b(-2+iy)=g_b(y^2),\qquad
 g_b(x)=-2\sum_{j=1}^b(-1)^{b-j}
                  \frac{(2j-1)/b}{x+((2j-1)/b)^2}.
\end{equation}
We will give these kernels a common positive denominator and use
their numerators as multiplier polynomials. This makes the constant
matrix an integral of products with one common weight.

\begin{example}\label{ex:two-kernels}
For $B=2$,
\[
 g_1(x)=-\frac2{x+1},\qquad
 g_2(x)=\frac1{x+1/4}-\frac3{x+9/4}.
\]
With $Q(x)=(x+1)(4x+1)(4x+9)$, write $g_b=h_b/Q$. Then
\[
 h_1(x)=-2(4x+1)(4x+9),\qquad
 h_2(x)=8(x+1)(3-4x).
\]
The four entries of the constant matrix will be integrals of
$h_1^2,h_1h_2,h_2h_1,h_2^2$ against the same positive weight.
This is the reason for choosing the kernel numerators as multipliers.
\end{example}

\subsection{Independent integer polynomials}

Let $\mathcal A$ be the set of distinct positive pole distances:
\[
 \mathcal A=\left\{\frac{2j-1}{b}:1\le b\le B,\ 1\le j\le b\right\},
 \qquad N=|\mathcal A|,\qquad J=N-1.
\]
For $a=v/d\in\mathcal A$, always use the reduced numerator and
denominator. Define
\begin{equation}\label{eq:Q-h}
 Q(x)=\prod_{v/d\in\mathcal A}(d^2x+v^2),\qquad
 h_b(x)=Q(x)g_b(x).
\end{equation}
Then $Q(x)>0$ for $x\ge0$, and $h_b\in\Z[x]$ has degree at most
$J$. Indeed, each summand of $g_b$ is of the form
$\pm2vd/(d^2x+v^2)$, whose denominator divides $Q$.

\begin{lemma}\label{lem:kernel-independence}
The polynomials $h_1,\ldots,h_B$ are linearly independent over $\R$.
Moreover,
\begin{equation}\label{eq:N-bound}
 N=D+\sum_{\substack{d\le B\\d\text{ even}}}\varphi(d)\le2D.
\end{equation}
\end{lemma}
\begin{proof}
In a nonzero relation among the $g_b$, choose the largest index $b$
with nonzero coefficient. The function $g_b$ has a pole at
$x=-1/b^2$. Every earlier $g_a$ is regular there, since all its
positive pole distances are at least $1/a>1/b$. This is impossible.
Multiplication by $Q$ preserves independence.

The elements of $\mathcal A$ are exactly the reduced fractions
$v/d$ with $d\le B$, $v$ odd and $0<v<2d$. There are $\varphi(d)$
such numerators for odd $d$, and $2\varphi(d)$ for even $d$.
Summing gives \eqref{eq:N-bound}.
\end{proof}

At the seed, $-(t+2)^2=y^2$ on the midpoint line. The multiplier
$(t+2)h_b(-(t+2)^2)$ therefore produces the desired factor $h_b(y^2)$
there, while its odd factor $t+2$ makes the product with $R_1$
reflection invariant. To lift this choice to every odd index, we
make it homogeneous in $n,t$:
\begin{align}
 W_{n,b}(t)
 &=n^{2J}(t+2n)h_b\!\left(-\frac{(t+2n)^2}{n^2}\right)\notag\\
 &=\sum_{k=0}^J h_{b,k}(-1)^k n^{2J-2k}(t+2n)^{2k+1},
 \qquad h_b(x)=\sum_{k=0}^J h_{b,k}x^k.\label{eq:W}
\end{align}
The expanded expression shows that the coefficients are integers.
Each term has the same total degree in $n,t$:
\begin{equation}\label{eq:E}
 E=2J+1=2N-1\le4D-1<s.
\end{equation}
The last inequality follows from $S>DL$ and $L\ge10$.
These polynomials are odd about $-2n$ and satisfy the useful scaling
identity
\begin{equation}\label{eq:homogeneity}
 W_{mp,b}(pt)=p^E W_{m,b}(t).
\end{equation}
The degree bound ensures decay at infinity; the scaling identity
will control prime dilation.

\subsection{The common positive weight}

Let $M_{b,j}$ be the $b$th rational constant associated with
$W_{1,j}R_1$. We now prove that the matrix $M=(M_{b,j})$ is
nonsingular.

Write $z=t+2$. The monic numerator $P_1$ has simple real roots and
is odd in $z$, so for a set $\Lambda$ of $DL/2$ positive rational
numbers,
\[
 P_1(t)=z\prod_{\alpha\in\Lambda}(z^2-\alpha^2).
\]
On the line $t=-2+iy$, the denominator $t^S(t+4)^S$ is
$-(y^2+4)^S$, since $S$ is odd. Substitution gives
\begin{equation}\label{eq:midpoint-value}
 W_{1,j}(-2+iy)R_1(-2+iy)
 =(-1)^{DL/2}C_1
  \frac{y^2\prod_{\alpha\in\Lambda}(y^2+\alpha^2)}{(y^2+4)^S}
  h_j(y^2).
\end{equation}

Every pole of $H_b$ is a point $-k/b\in(-4,0)$ where $P_1$ vanishes.
These poles cancel in $W_{1,j}R_1H_b$, leaving only $0$ and $-4$.
Close the upward midpoint line by a large semicircle to its right.
The contour encloses $0$ clockwise, so the residue theorem and
\eqref{eq:residue-constant} give
\begin{equation}\label{eq:Gram}
 M_{b,j}=\frac{\sigma C_1}{2\pi}
    \int_{-\infty}^{\infty}\omega(y)h_b(y^2)h_j(y^2)\,dy,
 \qquad \sigma=(-1)^{DL/2+1},
\end{equation}
where
\[
 \omega(y)=\frac{y^2\prod_{\alpha\in\Lambda}(y^2+\alpha^2)}
                  {(y^2+4)^S Q(y^2)}.
\]
For completeness, $H_b(t)=O(t^{-2})$ because its character weights
sum to zero. The contour integrand therefore has degree at most
$DL+2N-2S-2<-2$ at infinity. The arc integral vanishes, and the
integral in \eqref{eq:Gram} converges absolutely.

The weight $\omega(y)$ is positive for every real $y\ne0$.
For a nonzero real vector $v$, Lemma~\ref{lem:kernel-independence}
shows that $\sum_bv_bh_b$ is a nonzero polynomial. Consequently
\[
 \sigma v^{\mathsf T}Mv
 =\frac{C_1}{2\pi}\int_{-\infty}^{\infty}
      \omega(y)\left(\sum_bv_bh_b(y^2)\right)^2\,dy>0.
\]
Thus $\sigma M$ is positive definite and $M$ is nonsingular.

The entries of $M$ are rational, being constants of rational partial
fractions. We can therefore choose fixed integers $c_1,\ldots,c_B$
such that
\[
 M(c_1,\ldots,c_B)^{\mathsf T}=\eta e_B,\qquad \eta\in\Q\setminus\{0\}.
\]
For all positive odd $n$, set
\begin{equation}\label{eq:T}
 W_n=\sum_{j=1}^B c_jW_{n,j},\qquad T_n=W_nR_n.
\end{equation}
This gives the prescribed seed vector \eqref{eq:target}. The
coefficients $c_j$ choose one multiplier $W_n$ shared by all $B$
forms. The weights introduced in Section~\ref{sec:contradiction}
will combine those forms to cancel specified beta terms.
The $c_j$ remain fixed as $n$ varies. Their size can affect
fixed constants in estimates, but not the exponential decay rate.
